\documentclass[uplatex, a4paper, 11pt]{jsarticle}

% --- パッケージの読み込み ---
\usepackage[utf8]{inputenc}
\usepackage[dvipdfmx]{graphicx} % 画像（ChemDrawなど）の挿入用
\usepackage{amsmath, amssymb}   % 数式用
\usepackage[version=4]{mhchem}  % 化学反応式用（\ce{}で記述）
\usepackage{siunitx}            % 単位用（\SI{10.5}{g/mol}など）
\usepackage{geometry}           % 余白設定用
\geometry{left=25mm,right=25mm,top=30mm,bottom=30mm}

% --- タイトル情報 ---
\title{第○回 実験レポート：〇〇の合成と評価}
\author{福山 聖太}
\date{\today}

\begin{document}

\maketitle

\section{目的}
本実験では、……の合成を行い、その構造をNMRおよびIRにより特定することを目的とする。

\section{実験方法}
\subsection{試薬および器具}
特記ない限り、購入した試薬をそのまま用いた。
\begin{itemize}
    \item スチレン（精製済み）
    \item アゾビスイソブチロニトリル (AIBN)
\end{itemize}

\subsection{合成手順}
以下の反応式に従い、ラジカル重合を行った。
\begin{equation}
    % mhchemを使った重合反応の例
    \ce{\mathit{n} CH2=CH(Ph) ->[AIBN][\Delta] -[ CH2-CH(Ph) ]_{\mathit{n}}}
\end{equation}

\section{結果と考察}
得られた高分子の数平均分子量（$M_n$）は \SI{1.5e4}{g/mol}、分散度（$M_w/M_n$）は \SI{1.20}{} であった。

% 画像を挿入する場合のテンプレート（コメントアウト解除して使用）
% \begin{figure}[htbp]
%     \centering
%     \includegraphics[width=10cm]{scheme1.pdf}
%     \caption{合成スキーム}
%     \label{fig:scheme1}
% \end{figure}

\section{結論}
……を合成することに成功した。

\end{document}